\section{Extension to any fixed number of experts}
\label{sec:tournament}

The two-expert learner can be composed in a tournament. Its actual child
loss sequences are adaptive functions of the original observations and
need not be product-distributed. We instead introduce deterministic
modal paths for the analysis of replication, conditional only on tapes.

\begin{lemma}[A modal whole-path comparison]
\label{lem:modal-path}
Let $Z=f(X,R)$ take values in a finite alphabet, where the input $X$ has
a fixed law and the original tape $R$ is independent of $X$. For each
tape value $s$, define $p_s(z)=\PP_X\{f(X,s)=z\}$ and select a maximizer
$z^*(s)$ with a fixed lexicographic rule for ties. If $X,X'$ are
independent copies sharing $R$, then
\begin{equation}
 \PP_{R,X}\{f(X,R)\ne z^*(R)\}
 \le\PP_{R,X,X'}\{f(X,R)\ne f(X',R)\}.
 \label{eq:modal-path-bound}
\end{equation}
\end{lemma}
\begin{proof}
The probabilities $p_s(z)$ are measurable functions of $s$, by integration
of the measurable indicators defining $f$. There are finitely many
$z$, so lexicographic maximization gives a measurable selection.
For each fixed $s$, independence of $X,X'$ gives conditional
disagreement $1-\sum_z p_s(z)^2$. Since
\[
 \sum_z p_s(z)^2\le\max_z p_s(z),
\]
the single-input modal deviation $1-\max_zp_s(z)$ is no larger.
Integrating over the original tape proves
\eqref{eq:modal-path-bound}.
\end{proof}

The modal path may depend on the unknown input law. It is used only in
the analysis and is neither computed nor supplied to the learner.
In particular, \eqref{eq:modal-path-bound} is an averaged-tape
inequality, not a bound on deviation for every fixed tape.

\begin{theorem}[Tournament reduction]
\label{thm:tournament}
Suppose a two-expert learner, for every $r\in(0,1)$, has expected regret
at most $B(T,r)$ on every deterministic sequence of costs in $[0,1]^2$
and is $r$-replicable for every fixed time-varying product law on those
cost pairs. The expectation and replication probability use its entire
original input-independent tape.

For $n\ge2$, $\rho\in(0,1)$, and $T\ge2$, let
\[
 N=2^{\lceil\log_2 n\rceil},\qquad
 D_{\rm tr}=\log_2N,\qquad r=\frac{\rho}{N^2}.
\]
There is an $n$-expert learner with expected regret at most
$D_{\rm tr}B(T,r)$ on every deterministic cost sequence, whose entire
original-expert action path has disagreement probability at most
\begin{equation}
 \frac{r(N^2-1)}{3}<\rho
 \label{eq:tournament-replication-bound}
\end{equation}
under any fixed nonstationary product law on $[0,1]^{n\times T}$.
Coordinates within a round may be dependent.
\end{theorem}

\begin{proof}
Pad the original experts with duplicates to form $N$ leaves of a complete
balanced binary tree. Each leaf recommends its fixed original expert
index. Each internal node runs a copy of the two-expert learner with
parameter $r$ on its own independent subtape. At a round, a node selects
one of its two children and recommends the original index currently
recommended by that child. The root recommendation is played.

After the full original cost vector is revealed, every internal node
receives the costs of both current child recommendations and updates.
All nodes operate in full information at every round, including nodes
not selected by their ancestors. Consequently a child's recommendation
path depends only on past original costs and tapes in that child's
subtree; it depends on neither its parent's tape nor any ancestor's
choices. This also specifies a causal implementation: all current
recommendations are determined before the current cost vector is
revealed.

\paragraph{Regret.}
Fix a deterministic original cost sequence. Conditional on all tapes
strictly below an internal node, its two child recommendation paths and
their two cost sequences are deterministic and independent of its own
tape. The node's expected total cost is therefore at most the smaller
child total plus $B(T,r)$, and hence at most the total of either
particular child plus $B(T,r)$. Integrating the descendant tapes gives
\[
 \EE L_{\rm node}\le\EE L_{\rm child}+B(T,r)
\]
for either specified child. Iterating along the path to any fixed leaf
yields
\[
 \EE L_{\rm root}\le L_{\rm leaf}+D_{\rm tr}B(T,r).
\]
Choosing a leaf representing a best original expert proves the regret
bound. This argument keeps all tape averaging intact and makes no
assumption that any child regret is nonnegative.

\paragraph{Replication.}
Fix an arbitrary product law on the original cost vectors. For a parent,
let $q_L,q_R$ be bounds on the disagreement probabilities of its child
original-index paths under that law, using independent input streams
and the same subtree tapes. Condition only on all tapes strictly below
the parent. For each child, choose its canonical modal whole
original-index path under the fixed law and those fixed tapes, as in
Lemma~\ref{lem:modal-path}. These two paths are now deterministic
sequences of indices. Their conditional definition may use the entire
fixed law; it is an analysis device, not part of the causal algorithm.

Construct a shadow parent in each stream by evaluating the current
original cost vector at these two deterministic shadow indices and
running the parent learner on the resulting cost pairs with its original
parent tape. For every fixed choice of descendant tapes, the shadow
cost pairs are independent across time, since the original vectors are.
The two shadow streams are independent copies of the same product law.
The parent's own tape remains independent of these inputs and of the
fixed descendant tapes. Its two-expert guarantee therefore bounds the
shadow parent decision-path disagreement by $r$. Integrating over the
descendant tapes preserves this bound.

Lemma~\ref{lem:modal-path}, followed by a union bound over both streams
and both children, bounds the unconditional probability that any actual
child path differs from its shadow by
\[
 2(q_L+q_R).
\]
Only the averaged-tape shadow inequality is used here. We do not
condition on the complement of this deviation event, and we do not
assert that the actual adaptive child costs form a product law.

On the event of no child-shadow deviation, the actual and shadow parent
cost sequences coincide within each stream. Causality and use of the
same parent tape then imply that their parent decision paths coincide
within that stream. If the two shadow parent decision paths also agree,
the selected deterministic shadow child indices agree at every time,
so the actual parent original-index paths agree. This pointwise event
containment proves
\begin{equation}
 q_{\rm parent}\le r+2(q_L+q_R).
 \label{eq:tournament-recursion}
\end{equation}
No independence between left and right child paths is required.

Leaves have zero disagreement. In a complete balanced tree,
\eqref{eq:tournament-recursion} gives, at height $j$, the bound
$q_j\le r+4q_{j-1}$, with $q_0=0$. Hence
\[
 q_{\rm root}\le r\sum_{j=0}^{D_{\rm tr}-1}4^j
   =\frac{r(N^2-1)}{3}<\rho,
\]
as claimed. All probabilities include the entire original collection
of node tapes, including any rare branch coins used by the base
learners.
\end{proof}

Applying Theorem~\ref{thm:tournament} to the two-expert bound already
proved gives
\[
 \Reg_T=O_{n,\rho}\!\left(\sqrt T\log^{7/2}(2T)\right)
\]
for every fixed finite $n$ and fixed $\rho\in(0,1)$. The dependence on
$n$ and $\rho$ is poor; the reduction makes no claim of an optimal
dependence on either parameter. For $n=1$, always select the sole expert.
For $T=1$, any shared input-independent expert choice replicates exactly
and has regret at most one. These conventions cover the finite boundary
cases without the tournament parameters.
